\section{Свойства преобразования Лапласа: подобие, сдвиг изображения,
запаздывание
оригинала.}\label{ux441ux432ux43eux439ux441ux442ux432ux430-ux43fux440ux435ux43eux431ux440ux430ux437ux43eux432ux430ux43dux438ux44f-ux43bux430ux43fux43bux430ux441ux430-ux43fux43eux434ux43eux431ux438ux435-ux441ux434ux432ux438ux433-ux438ux437ux43eux431ux440ux430ux436ux435ux43dux438ux44f-ux437ux430ux43fux430ux437ux434ux44bux432ux430ux43dux438ux435-ux43eux440ux438ux433ux438ux43dux430ux43bux430.}

\textbf{Подобие:} Если \(f(t)\risingdotseq F(p)\) и \(\beta>0\), то \[
f(\beta t)\risingdotseq \frac{1}{\beta}F\left( \frac{p}{\beta} \right)
\] \textbf{Сдвиг изображения:} Если \(f(t)\risingdotseq F(p)\), то
\(\forall a\in \mathbb{C}:\) \[
f(t)e^{ a t }\risingdotseq F(p-a)
\]

\textbf{Запаздывание оригинала:} Если \(f(t)\risingdotseq F(p)\) и
\(f(t)=0\) при \(t<\tau\), где \(\tau>0\), то \[
f(t-\tau)\risingdotseq e^{ -p\tau }F(p).
\]
